\section{Evaluation}
\label{sec:evaluation}

\subsection{Test Design}
\label{sec:test-design}

The evaluation is a small internal pilot, not a user study. Two reviewers each processed 48
documents: half through the manual workflow first, half through the assisted workflow first,
with the order randomised to reduce learning effects. We recorded three measures: median
review time, the share of model proposals overridden by the reviewer, and evidence
completeness (the share of final labels for which the review record contains a supporting
text span). All documents are synthetic; the evaluation involves no real individuals.

\subsection{Results}
\label{sec:results}

Table~\ref{tab:results} summarises the results. The assistant reduced the median review time
from 6.5 to 3.8 minutes and raised evidence completeness from 62\% to 91\%. The cost was 9
overrides out of 48 items, an override rate of 18.8\%. That number should not be read as
``the model is 81.2\% accurate'': an override happens whenever the model and the reviewer
disagree, and the disagreement can be a model error or a reviewer changing their judgement.

\begin{table}[htbp]
  \centering
  \begin{tabular}{lccc}
    \toprule
    Metric & Manual & Assisted & Change \\
    \midrule
    Median review time (min) & 6.5 & 3.8 & $-42\%$ \\
    Reviewer override rate & -- & 18.8\% & -- \\
    Evidence completeness & 62\% & 91\% & $+29$ pp \\
    \bottomrule
  \end{tabular}
  \caption{Results from the small internal pilot. The sample is 48 synthetic documents
  reviewed by two reviewers.}
  \label{tab:results}
\end{table}

\subsection{Failure Cases}
\label{sec:failures}

The failures cluster into three groups:
\begin{itemize}
  \item \textbf{Mixed topics.} When a document covers two topics, the model emits one label
        and the reviewer has to split it;
  \item \textbf{Missing evidence.} A very short document receives a high score but has no
        usable text span;
  \item \textbf{Ambiguous labels.} The boundary between two topics is genuinely unclear, so
        the disagreement cannot be attributed to a model error.
\end{itemize}

Table~\ref{tab:test-cases} in the appendix lists the individual test records. We keep these
failure cases rather than reporting only the successful sample, because a reviewer needs to
know when the system is unreliable.
